
% Tabela em arquivo separado. Não é um documento autônomo.

% Base visual obrigatória: ../../03_eixos_tematicos.tex.

% Fontes examinadas: meu-trabalho.pdf (páginas físicas 11--37),

% elementos-textuais/fundamentacao-teorica.tex, trabalhos-relacionados.tex,

% metodologia.tex, referencias.bib, meu-trabalho.tex e cefetmg.cls.

% A Tabela 1 de Maria Rocha (PDF, pp. 28--30) orienta apenas as cinco colunas.

%

% DEPENDÊNCIAS (nenhum preâmbulo foi alterado por esta revisão):

% Já carregados no projeto: booktabs, color, colortbl (que carrega array),

% caption, abntex2cite; fonte Helvetica e recursos da classe cefetmg/abntex2.

% Necessários: longtable, ragged2e, booktabs e colortbl. Na conferência final,

% estes pacotes já constavam do preâmbulo, por alteração externa à revisão.

% Todos estavam instalados no ambiente de validação; não há nova dependência.

% Usar este conteúdo no corpo do documento, fora de floats. A inclusão

% no capítulo 3 foi encontrada durante a conferência; não foi criada aqui.

% Não envolver em table, minipage ou resizebox. Referência: \ref{tab:sintese_referencial}.

%

% ESCOPO: 32 referências únicas; seleção temática, não exaustiva.

% A síntese registra o papel atribuído às obras nas seções 2 e 3 do TCC;

% não constitui nova verificação integral de todos os artigos originais.

% Revisão: as sete inclusões foram cotejadas com o TCC e os PDFs locais

% dos artigos. Ma, Wang e Brauneis também tiveram pontos conferidos no original.

% Cui (2023, seção 3.1 do artigo): PPO, não "modelo proprietário";

% Mean-CVaR é comparado à MV, não se trata de fracasso geral de ML + MV.

% Zouaoui: LSTM estima pesos; não foi descrita como simples previsão de retornos.

% Mba: os resultados negativos se limitam à amostra e às regras testadas.

% Ma: ARIMA supera CNN e SVR na previsão; após custos, a preferência por

% Random Forest em MV não se estende automaticamente ao modelo ômega.

% As sete candidatas foram incluídas por acrescentarem aspectos distintos;

% as demais omissões mantêm o caráter seletivo, não exaustivo, da síntese.

% Trechos em comment e lembretes todo não foram tratados como evidência.

% Frydman (1985) e White (1988) são antecedentes mencionados por revisões,

% não entradas diretamente consultadas ou novas chaves bibliográficas.

% A tabela ilustrativa da versão inicial não foi usada por ausência de fonte.

% A metodologia prevê AdaBoost + MV versus 1/N e teste fora da amostra;

% não se atribuem ao TCC testes de custos, risco de cauda ou arquiteturas extras.

% Ma (2021): a seção 2 ressalta ARIMA e custos; a seção 3 destaca Random Forest.

% Preservou-se essa nuance, sem afirmar superioridade universal de ML.

% Assaf Neto: o estilo atual renderiza a entrada como "Neto (2014)".

% Abaixo, o sobrenome composto é explicitado localmente, com \citeyear,

% preservando integralmente o .bib e as configurações bibliográficas.

%

% DIAGRAMAÇÃO: eixos e paleta preservados; tabcolsep=3.1pt, arraystretch=1.3,

% réguas de 1.6/0.5/0.4/1pt e faixa temática com 1pt seguem o modelo.

% Adaptações: fonte 9pt, margens mantidas, nenhuma geometria local,

% retrato em A4, sem alteracao de geometria; colunas usando largura disponivel.

% Mantém-se os 32 estudos. A origem de cada obra aparece abaixo da autoria; referências comuns às

% duas seções são reunidas em uma única linha. Subdivisões refinam o eixo 3.

\begingroup

\SingleSpacing

\fontsize{9}{11}\selectfont

\setlength{\parindent}{0pt}

\setlength{\parskip}{0pt}

\setlength{\tabcolsep}{3.1pt}

\renewcommand{\arraystretch}{1.3}

\setlength{\LTpre}{0pt}

\setlength{\LTpost}{0pt}

\setlength{\LTleft}{0pt plus 1fill}

\setlength{\LTright}{0pt plus 1fill}

\setlength{\LTcapwidth}{\linewidth}

\setlength{\emergencystretch}{3em}

\definecolor{sinteseForest}{rgb}{0.117647,0.286275,0.262745}

\definecolor{sinteseSoft}{rgb}{0.909804,0.941176,0.933333}

\definecolor{sinteseLine}{rgb}{0.784314,0.847059,0.831373}

\newcolumntype{S}[1]{>{\RaggedRight\arraybackslash}p{#1}}

% As proporções mantêm a hierarquia do modelo e somam exatamente 1.

\def\SinteseLargura{\dimexpr\linewidth-8\tabcolsep\relax}

\newcommand{\SinteseAutor}[2]{\textbf{\citeonline{#1}}\par

   {\normalfont\itshape Seções #2}}

\newcommand{\SinteseEixo}[1]{%

   \rowcolor{sinteseSoft}\multicolumn{5}{@{}l}{%

   \color{sinteseForest}\bfseries #1}\\*[1pt]}

\newcommand{\SinteseLinha}{\arrayrulecolor{sinteseLine}\specialrule{0.4pt}{0pt}{0pt}}

\newcommand{\SinteseCabecalho}{%

   \color{sinteseForest}\bfseries Autor(es) &

   \color{sinteseForest}\bfseries Foco do estudo &

   \color{sinteseForest}\bfseries Método / \newline contexto &

   \color{sinteseForest}\bfseries Principais \newline contribuições &

   \color{sinteseForest}\bfseries Uso no presente estudo\\}

\begin{longtable}{
    @{}
    S{0.175\SinteseLargura}
    S{0.125\SinteseLargura}
    S{0.2\SinteseLargura}
    S{0.26\SinteseLargura}
    S{0.26\SinteseLargura}
    @{}
}

\caption{Síntese temática da fundamentação e dos estudos relacionados}\label{tab:sintese_referencial}\\

\arrayrulecolor{sinteseForest}\toprule[1.6pt]

\SinteseCabecalho

\arrayrulecolor{sinteseLine}\midrule[0.5pt]

\endfirsthead

\multicolumn{5}{@{}l}{\itshape Tabela \thetable}\\[1mm]

\arrayrulecolor{sinteseForest}\toprule[1.6pt]

\SinteseCabecalho

\arrayrulecolor{sinteseLine}\midrule[0.5pt]

\endhead

\endfoot

\arrayrulecolor{sinteseForest}\bottomrule[1pt]
\multicolumn{5}{@{}p{\linewidth}@{}}{%
    \itshape Fonte: elaboração própria.}\\
\endlastfoot

\SinteseEixo{FUNDAMENTOS DE FINANÇAS}

% markowitz1952portfolio | seções 2.1; 2.4

\begin{comment}


\SinteseAutor{markowitz1952portfolio}{2.1; 2.4} &

Seleção de carteiras sob risco. &

Formulação teórica de média-variância (MV); retorno esperado, variância e covariância. &

Estabelece a alocação conjunta de ativos a partir do compromisso entre risco e retorno. &

Fundamenta o componente MV da estratégia AdaBoost + MV proposta no TCC. \\


\SinteseLinha
\end{comment}


% assaf2014financas | seções 2.4; 2.4.1

\textbf{Assaf Neto (\citeyear{assaf2014financas})}\par {\normalfont\itshape Seções 2.4; 2.4.1} &

Risco, diversificação e fronteira eficiente. &

Exposição teórica de finanças; correlação entre ativos e curvas de indiferença. &

Explica a composição do risco da carteira, os limites da diversificação e a escolha entre carteiras eficientes. &

Sustenta a formulação do risco e a interpretação dos gráficos de combinações de ativos e da fronteira eficiente. \\

\SinteseLinha

% ross2006fundamentals | seções 2.4

\SinteseAutor{ross2006fundamentals}{2.4} &

Componentes do risco financeiro. &

Fundamentos de finanças corporativas; risco sistemático e não sistemático. &

Distingue choques comuns ao mercado de riscos específicos de ativos ou grupos de ativos. &

Delimita o risco que a diversificação pode mitigar; impede interpretar a otimização como eliminação de todo risco. \\

\SinteseLinha

% almansour2023behavioral | seções 2.1

\SinteseAutor{almansour2023behavioral}{2.1} &

Vieses e decisões de investimento. &

Finanças comportamentais; influências sociais e percepção de risco. &

Relaciona fatores psicológicos e mudanças na percepção de risco a decisões que se afastam da racionalidade clássica. &

Contextualiza os limites da hipótese de investidor plenamente racional na discussão sobre tomada de decisão. \\

\SinteseLinha

% charfeddine2020investigating | seções 2.1; 2.2

\SinteseAutor{charfeddine2020investigating}{2.1; 2.2} &

Criptoativos e ativos convencionais. &

Retornos diários de 2015 a 2018; Bitcoin, Ethereum e mercados tradicionais. &

Documenta elevada volatilidade e baixa correlação em parte das combinações, com potencial de diversificação. &

Caracteriza o mercado investigado e sustenta a discussão sobre dependência entre ativos e composição de carteiras. \\

\SinteseLinha

% goodell2021diversifying | seções 2.2

\SinteseAutor{goodell2021diversifying}{2.2} &

Diversificação durante crises. &

Criptomoedas, stablecoins e índices acionários durante a COVID-19. &

Relata aumento de comovimento entre criptomoedas convencionais e ações. &

Qualifica a expectativa de diversificação: os benefícios dependem do tipo de criptoativo e das condições de mercado. \\

\SinteseLinha

\SinteseEixo{ML APLICADO A FINANÇAS}

% aranha2023efficacies | seções 2.3.3

\SinteseAutor{aranha2023efficacies}{2.3.3} &

Antecedentes do ML financeiro. &

Revisão de modelos de risco de crédito e previsão de falência. &

Retoma o particionamento recursivo de Frydman et al. (1985) como antecedente das árvores de decisão em finanças. &

Situa historicamente o uso dos métodos baseados em árvores em finanças. \\

\SinteseLinha

% balasubramanian2024survey | seções 2.3.3

\SinteseAutor{balasubramanian2024survey}{2.3.3} &

Evolução da previsão acionária. &

Revisão da literatura sobre previsão do mercado de ações. &

Relata a experiência pioneira de White (1988) com redes neurais. &

Contextualiza a evolução do ML e seus limites iniciais. \\

\SinteseLinha

% busari2021crude | seções 2.3.2

\SinteseAutor{busari2021crude}{2.3.2} &

Combinação adaptativa de preditores. &

AdaBoost-LSTM e AdaBoost-GRU na previsão do preço do petróleo. &

Oferece a base usada no texto para explicar treinamento sequencial e reponderação de observações no AdaBoost. &

Fundamenta o funcionamento do algoritmo escolhido para a etapa preditiva. \\

\SinteseLinha

% park2023intelligent | seções 2.2; 2.3; 3.2

\SinteseAutor{park2023intelligent}{2.2; 2.3; 3.2} &

Previsão direcional e negociação. &

AdaBoost-LSTM com informação de regimes de turbulência; negociação de um único criptoativo. &

Relata ganhos frente a indicadores técnicos em parte dos horizontes; discute a sensibilidade do AdaBoost a ruído e valores atípicos. &

Oferece antecedente do AdaBoost em criptoativos e fundamenta suas limitações; a extensão para alocação entre ativos pertence à proposta do TCC. \\

\SinteseLinha

% derbentsev2021comparative | seções 2.3; 3.2

\SinteseAutor{derbentsev2021comparative}{2.3; 3.2} &

Previsão com conjuntos de árvores. &

Random Forest e boosting; preços diários de BTC, ETH e XRP; teste de um passo à frente. &

Apresenta erros preditivos reduzidos e ganhos com atributos adicionais, sem avaliar rentabilidade de carteiras. &

Sustenta os conceitos de ensemble e avaliação fora da amostra; delimita a diferença entre precisão de preços e qualidade da alocação. \\

\SinteseLinha

% livieris2020ensemble | seções 2.1; 2.3.1; 3.2

\SinteseAutor{livieris2020ensemble}{2.1; 2.3.1; 3.2} &

Previsão por conjuntos de redes neurais. &

CNN, LSTM e BiLSTM combinadas por média, bagging e stacking; preços horários de BTC, ETH e XRP. &

Relata ganhos preditivos com maior custo de treinamento e previsão; não demonstra ganhos de rentabilidade. &

Fundamenta a combinação de preditores e explicita o custo computacional; reforça a avaliação da carteira, além do erro de previsão. \\

\SinteseLinha

% sebastiao2021forecasting | seções 2.2; 2.3.1; 3.2

\SinteseAutor{sebastiao2021forecasting}{2.2; 2.3.1; 3.2} &

Previsão sob mudanças de mercado. &

Modelos lineares, Random Forest e SVM; BTC, ETH e LTC; janelas móveis. &

Não identifica superioridade universal; custos de negociação tornam negativas algumas estratégias antes rentáveis. &

Fundamenta a separação temporal de treino e teste e a avaliação econômica dos sinais, além da precisão preditiva. \\

\SinteseLinha

% chevallier2021forecast | seções 2.2; 3.2

\SinteseAutor{chevallier2021forecast}{2.2; 3.2} &

Previsão e negociação de Bitcoin. &

Comparação de modelos, incluindo AdaBoost; informações de outras criptomoedas e classes de ativos. &

Ganhos de previsão não asseguram ganhos de negociação; buy-and-hold supera, em geral, as regras testadas. &

Reforça a necessidade de avaliar o AdaBoost pelo desempenho da estratégia resultante, sem presumir vantagem econômica por menor erro. \\

\SinteseLinha

\SinteseEixo{ESTUDOS RELACIONADOS --- PREVISÃO E ALOCAÇÃO}

% wang2020portfolio | seções 2.1; 2.4; 3.1

\SinteseAutor{wang2020portfolio}{2.1; 2.4; 3.1} &

Pré-seleção e formação de carteiras. &

LSTM aplicada a séries financeiras de longo prazo, seguida de seleção de ativos e otimização MV. &

Relata melhorias em retorno cumulativo e ajustado ao risco na maior parte dos períodos analisados. &

Exemplifica a sequência previsão, seleção e alocação que aproxima o ML da estrutura de Markowitz. \\

\SinteseLinha

% ma2021portfolio | seções 2.3.3; 3.1

\SinteseAutor{ma2021portfolio}{2.3.3; 3.1} &

Previsão de retornos e otimização. &

Comparação de modelos estatísticos, ML e aprendizado profundo; seleção e alocação por MV e ômega. &

ARIMA supera CNN e SVR na previsão; os ganhos de Random Forest dependem da alocação e dos custos, que favorecem SVR em ômega. &

Articula as duas seções ao distinguir qualidade preditiva de utilidade na carteira; evita pressupor superioridade automática do ML. \\

\SinteseLinha

% du2022mean | seções 2.1; 3.1

\SinteseAutor{du2022mean}{2.1; 3.1} &

Sensibilidade da alocação às estimativas. &

Redes LSTM com atenção para previsão de portfólios de ações cointegradas, antes da otimização MV. &

Discute a sensibilidade da MV a erros nas médias e covariâncias e propõe aprimorar suas entradas por previsão. &

Explicita o problema de estimação que motiva combinar um preditor com a otimização clássica. \\

\SinteseLinha

% chaweewanchon2022markowitz | seções 2.1; 2.4; 3.1

\SinteseAutor{chaweewanchon2022markowitz}{2.1; 2.4; 3.1} &

Pré-seleção preditiva antes da MV. &

CNN e BiLSTM para selecionar ações antes da alocação média-variância. &

Relata melhorias em Sharpe, retorno médio e risco nos testes; separa a seleção dos ativos da otimização dos pesos. &

Oferece antecedente de integração por pré-seleção e sustenta a discussão dos limites práticos dos pesos produzidos pela MV. \\

\SinteseLinha

% slusarczyk2025optimal | seções 3.1

\SinteseAutor{slusarczyk2025optimal}{3.1} &

Retornos históricos versus previstos. &

ARIMA-GARCH e XGBoost; carteiras de máximo índice de informação e mínima variância global. &

Retornos previstos podem melhorar a alocação, mas o ganho depende do controle de parâmetros e hiperparâmetros. &

Qualifica a expectativa de ganho ao substituir estimativas históricas por previsões; não assegura vantagem do AdaBoost no mercado cripto. \\

\SinteseLinha

% jiang2024deep | seções 2.4.1; 3.1

\SinteseAutor{jiang2024deep}{2.4.1; 3.1} &

Decisão dinâmica de portfólio. &

Aprendizado por reforço profundo; recompensa com extensão da MV, aversão ao risco e custos de transação. &

Integra risco e custos à aprendizagem da alocação e distingue risco da carteira de aversão individual ao risco. &

Apoia a discussão de preferências do investidor e oferece contraponto ao fluxo em duas etapas de AdaBoost + MV. \\

\SinteseLinha

\SinteseEixo{ESTUDOS RELACIONADOS --- MV NO MERCADO DE CRIPTOATIVOS}

% brauneis2019cryptocurrency | seções 2.2; 3.3

\SinteseAutor{brauneis2019cryptocurrency}{2.2; 3.3} &

Diversificação e desempenho da MV. &

Universo de 500 criptomoedas entre 2015 e 2017; comparação com ativos isolados e carteira 1/N. &

Fora da amostra, a carteira 1/N supera mais de 75\% das carteiras MV avaliadas em Sharpe e equivalente de certeza. &

Oferece justificativa empírica para comparar AdaBoost + MV com 1/N, referência efetivamente prevista na metodologia. \\

\SinteseLinha

% jeleskovic2024cryptocurrency | seções 2.4; 3.3

\SinteseAutor{jeleskovic2024cryptocurrency}{2.4; 3.3} &

Dependência e otimização em criptoativos. &

Modelo GARCH-cópula no arcabouço de Markowitz; carteiras com criptoativos e ativos tradicionais. &

Aponta benefícios em carteiras mistas e dificuldades para estimar dependências em um mercado de correlações instáveis. &

Sustenta a exposição sobre MV e Sharpe e contextualiza as limitações de estimar o risco das carteiras de criptoativos. \\

\SinteseLinha

% zouaoui2025portfolio | seções 3.3

\SinteseAutor{zouaoui2025portfolio}{3.3} &

Concentração e limites de Markowitz. &

Comparação entre Markowitz e abordagem com LSTM para estimar pesos de carteiras de criptomoedas. &

Relata concentração em poucos ativos e Sharpe menos favorável na carteira Markowitz avaliada; a LSTM obtém melhor Sharpe no experimento. &

Acrescenta a concentração da carteira às limitações da MV discutidas no TCC; o resultado com LSTM não valida diretamente AdaBoost + MV. \\

\SinteseLinha

\SinteseEixo{ESTUDOS RELACIONADOS --- MODELOS HÍBRIDOS EM CRIPTOATIVOS}

% zhou2023multisource | seções 3.2; 3.4

\SinteseAutor{zhou2023multisource}{3.2; 3.4} &

Informação preditiva na alocação. &

SVM com volume e sentimento; sete criptoativos; variância mínima modificada; testes de 1, 3 e 5 dias. &

Relata melhor desempenho em Sharpe, Sortino e equivalente de certeza frente às referências avaliadas, incluindo 1/N. &

Oferece antecedente direto da integração previsão + MV em criptoativos e de sua comparação com pesos iguais. \\

\SinteseLinha

% lorenzo2023online | seções 2.4.1; 3.2; 3.4

\SinteseAutor{lorenzo2023online}{2.4.1; 3.2; 3.4} &

Seleção de subconjuntos de ativos. &

Agrupamento por protótipos antes da MV; universos de 175, 250 e 534 criptoativos elegíveis. &

Reduz a dimensão da matriz de covariância e falhas de convergência, mas pode elevar perdas máximas e restringir a diversificação. &

Diferencia o uso de ML para seleção de ativos de seu uso preditivo; também apoia a interpretação do índice de Sharpe. \\

\SinteseLinha

% bulani2025improving | seções 3.4

\SinteseAutor{bulani2025improving}{3.4} &

Seleção por agrupamentos e alocação. &

K-medoids para agrupar criptoativos; seleção entre grupos e variantes de otimização por enxame de partículas. &

A seleção dinâmica entre grupos melhora Sharpe e Sharpe ajustado nos testes; faltam testes formais de significância. &

Distingue agrupamento e busca dos pesos da previsão de retornos; oferece alternativa metodológica para discutir diversificação e robustez. \\

\SinteseLinha

% toscano2026richer | seções 3.4

\SinteseAutor{toscano2026richer}{3.4} &

Valor de informações adicionais. &

ML com MV e MV modificada em tokens do ecossistema NFT; comparação com 1/N e custos. &

Ambas as estruturas superam 1/N nos testes descritos; a MV com ML, sem a modificação, geralmente apresenta melhor Sharpe e menor risco. &

Oferece comparação híbrida com 1/N e mostra que ampliar a informação utilizada não garante melhor alocação. \\

\SinteseLinha

% peykani2026machine | seções 3.4

\SinteseAutor{peykani2026machine}{3.4} &

Previsão e diferentes medidas de risco. &

XGBoost com ajuste bayesiano; três criptomoedas e duas ações; MV e outras duas formulações de risco. &

No teste de 30 dias, o retorno acumulado é 7,2\% para máximo Sharpe e 6,6\% para 1/N; os resultados variam com a medida de risco. &

Fornece antecedente da comparação entre previsão + otimização e 1/N; a carteira mista e o teste curto limitam a transposição dos achados. \\

\SinteseLinha

% xu2025cryptocurrency | seções 3.2; 3.4

\SinteseAutor{xu2025cryptocurrency}{3.2; 3.4} &

Previsão recorrente e Markowitz. &

LSTM; BTC, ETH e LTC; janela de 128 observações e pesos mensais com possibilidade de venda a descoberto. &

Relata retorno ligeiramente maior, risco semelhante e Sharpe superior em quatro dos seis meses avaliados. &

Introduz evidência de ganhos modestos em uma arquitetura próxima; impede tratar a combinação ML + MV como melhoria assegurada. \\

\SinteseLinha

% mba2026regime | seções 3.4

\SinteseAutor{mba2026regime}{3.4} &

Limites da alocação com redes em grafos. &

Redes neurais em grafos de correlação integradas à alocação; 33 pares cripto/USDT da Binance, de 2021 a 2026. &

Na amostra, mínima variância e ponderação por capitalização superam as configurações em grafos nas principais métricas de risco e retorno. &

Oferece evidência desfavorável à integração ML + MV e reforça a comparação com regras simples, sem presumir ganhos pela complexidade. \\

\SinteseLinha

% cui2023portfolio | seções 3.2; 3.4

\SinteseAutor{cui2023portfolio}{3.2; 3.4} &

Escolha da medida de risco. &

Aprendizado por reforço profundo; seis criptoativos; comparação \acs{mv} e \textit{Mean}-\acs{cvar}. &

\textit{Mean}-\acs{cvar} apresenta maiores retornos médios e menores desvios-padrão que \acs{mv} nas carteiras testadas. &

Expõe a limitação de representar risco apenas pela variância; a comparação favorece outra medida de risco. \\

\SinteseLinha

% chen2025multitask | seções 3.4

\SinteseAutor{chen2025multitask}{3.4} &

Aprendizagem dos pesos da carteira. &

LSTM para previsão e rede neural multitarefa para produzir os pesos; comparação com MV tradicional. &

Relata ganhos em retorno e Sharpe, acompanhados de resultados desfavoráveis em algumas métricas de risco. &

Distingue a produção de pesos por rede neural da proposta de usar AdaBoost na previsão e MV na alocação. \\

\end{longtable}

\endgroup


